% ============================================================
% CAPITOLO A - FONDAMENTI
% ============================================================
\chapter{Fondamenti}
\label{chap:fondamenti}

L'armonia tonale, prima ancora di essere un insieme di regole stilistiche, è un sistema logico basato su fenomeni acustici naturali. Per comprendere la sua grammatica, è indispensabile partire dalla fisica del suono e dalle relazioni matematiche che legano le note tra loro. Questo capitolo esplora i mattoni fondamentali del nostro sistema musicale: gli armonici, gli intervalli e la struttura gerarchica della scala.

\section{La natura del suono}
\label{sec:natura_suono}

Ogni suono che percepiamo come "musicale", ovvero di altezza determinata, non è un'entità singola, ma un fenomeno complesso. Quando una corda di chitarra o la colonna d'aria in un flauto vibrano, non lo fanno solo per la loro intera lunghezza, ma anche per frazioni intere di essa (metà, un terzo, un quarto, e così via). Ognuna di queste vibrazioni secondarie produce un suono più acuto e meno intenso, che si fonde con il suono principale.

\begin{definizione}[La serie degli armonici]
	\label{def:serie_armonici}
	La \textbf{serie degli armonici naturali} è la successione di suoni secondari che accompagnano un suono fondamentale. Questi suoni, detti \textbf{armonici} (o ipertoni), hanno frequenze (percepite dal nostro orecchio come altezza del suono) che sono multipli interi della frequenza del suono fondamentale.
\end{definizione}

L'ordine in cui gli armonici appaiono è universale e immutabile. Se prendiamo come fondamentale un Do, i primi suoni che emergono dalla sua vibrazione sono, in ordine: un Do più acuto, un Sol, un altro Do, un Mi, un altro Sol e un \Sibemolle  (leggermente calante rispetto al sistema musicale moderno).

\begin{esempio}[I primi 7 suoni armonici di Do]
	\label{ex:serie_armonici_do}
	% Questo esempio non è SATB, quindi usa la sintassi standard di LilyPond
	\musica{es_cap01_01}
\end{esempio}

Come si può osservare nell'Esempio \ref{ex:serie_armonici_do}, la natura stessa ci fornisce gli intervalli fondamentali dell'armonia. Dopo l'ottava (il raddoppio della fondamentale), il primo nuovo suono a comparire è la \textbf{quinta giusta} (Do-Sol), seguito dalla \textbf{terza maggiore} (Do-Mi). Questi tre suoni (Do, Mi, Sol), presi insieme, formano l'accordo perfetto maggiore (come scopriremo nel Capitolo \ref{chap:triadi}), la cellula base di tutta la musica tonale. La serie armonica, quindi, non è un'astrazione teorica, ma la giustificazione fisica dell'esistenza e della stabilità dell'accordo maggiore.

È interessante notare che il settimo armonico è un \Sibemolle  "stonato" rispetto al nostro sistema. Questo ci introduce a un concetto importante: il \textbf{temperamento}. Per poter suonare in tutte le tonalità su strumenti a intonazione fissa come il pianoforte o la chitarra, è stato necessario "temperare" gli intervalli, ovvero alterarli leggermente rispetto ai loro rapporti fisici puri. Il sistema oggi in uso, il \textbf{temperamento equabile}, divide l'ottava in dodici semitoni perfettamente uguali, sacrificando la purezza acustica in favore della flessibilità.

Più avanti analizzeremo la sequenza degli armonici prodotti da alcuni particolari gradi della scala musicale, per cui ne diamo adesso un elenco:

\begin{itemize}
	\label{itemize:elenco_armonici}
	
	\item \textbf{Do}: Do, Do, Sol, Do, Mi, Sol, La/\Ladiesis , Do
	
	\item \textbf{\Dodiesis }: \Dodiesis , \Dodiesis , \Soldiesis , \Dodiesis , Fa, \Soldiesis , \Ladiesis /Si, \Dodiesis 
	
	\item \textbf{Re}: Re, Re, La, Re, \Fadiesis , La, Si/Do, Re
	
	\item \textbf{\Rediesis }: \Rediesis , \Rediesis , \Ladiesis , \Rediesis , Sol, \Ladiesis , Do/\Dodiesis , \Rediesis 
	
	\item \textbf{Mi}: Mi, Mi, Si, Mi, \Soldiesis , Si, \Dodiesis /Re, Mi
	
	\item \textbf{Fa}: Fa, Fa, Do, Fa, La, Do, Re/\Rediesis , Fa
	
	\item \textbf{\Fadiesis }: \Fadiesis , \Fadiesis , \Dodiesis , \Fadiesis , \Ladiesis , \Dodiesis , \Rediesis /Mi, \Fadiesis 
	
	\item \textbf{Sol}: Sol, Sol, Re, Sol, Si, Re, Mi/Fa, Sol
	
	\item \textbf{\Soldiesis }: \Soldiesis , \Soldiesis , \Rediesis , \Soldiesis , Do, \Rediesis , Fa/\Fadiesis , \Soldiesis 
	
	\item \textbf{La}: La, La, Mi, La, \Dodiesis , Mi, \Fadiesis /Sol, La
	
	\item \textbf{\Ladiesis }: \Ladiesis , \Ladiesis , Fa, \Ladiesis , Re, Fa, Sol/\Soldiesis , \Ladiesis 
	
	\item \textbf{Si}: Si, Si, \Fadiesis , Si, \Rediesis , \Fadiesis , \Soldiesis /La, Si
	
\end{itemize}


\section{Gli intervalli}
\label{sec:intervalli}

Un intervallo è la distanza tra due suoni. Per misurarlo, usiamo due criteri: uno quantitativo (il numero) e uno qualitativo (la specie).

\begin{definizione}[Nomenclatura degli intervalli]
	\label{def:intervalli}
	Un intervallo si definisce contando il numero di gradi della scala compresi tra le due note (estremi inclusi). La sua specie può essere:
	\begin{itemize}
		\item \textbf{Maggiore (M)} o \textbf{Minore (m)}: per intervalli di seconda, terza, sesta e settima.
		\item \textbf{Giusto (G)}: per intervalli di unisono, quarta, quinta e ottava.
		\item \textbf{Aumentato (Eccedente)} o \textbf{Diminuito}: quando un intervallo maggiore, minore o giusto viene ampliato o ristretto cromaticamente di un semitono.
	\end{itemize}
\end{definizione}

Gli intervalli si classificano anche in base alla loro sonorità: le \textbf{consonanze} producono un effetto di stabilità e riposo, mentre le \textbf{dissonanze} generano tensione e un bisogno di risoluzione.

\subsection{I rivolti e la regola del 9}

Ogni intervallo può essere trasformato invertendo la posizione delle sue due note. Questa operazione, fondamentale sia nella teoria che nella pratica musicale, prende il nome di \textbf{rivolto} (o inversione).

\begin{definizione}[Rivolto di un intervallo]
	\label{def:rivolto_intervallo}
	Il \textbf{rivolto} (o \textbf{inversione}) di un intervallo si ottiene spostando la nota più grave un'ottava sopra (o equivalentemente, la nota più acuta un'ottava sotto). In questo modo, la nota che era al grave diventa acuta e viceversa.
\end{definizione}

Il rivolto non è una semplice trasposizione: cambia la natura stessa dell'intervallo. Una terza maggiore, ad esempio, diventa una sesta minore; una quinta giusta diventa una quarta giusta. Esiste tuttavia una relazione aritmetica precisa che governa questa trasformazione.

\begin{regola}[Aritmetica dei rivolti: la Regola del 9]
	\label{reg:regola_del_9}
	La somma di un intervallo e del suo rivolto è sempre 9. La qualità si trasforma secondo il seguente schema:
	\begin{itemize}
		\item \textbf{Maggiore} diventa \textbf{Minore} (e viceversa).
		\item \textbf{Giusto} rimane \textbf{Giusto}.
		\item \textbf{Aumentato} diventa \textbf{Diminuito} (e viceversa).
	\end{itemize}
\end{regola}

\begin{esempio}[Inversione di una Terza Maggiore]
	\label{ex:inversione_intervallo}
	Una terza maggiore (Do-Mi) si rivolta in una sesta minore (Mi-Do). Aritmeticamente: $3+6=9$.
	\musica{es_cap01_02}
\end{esempio}

Questo principio è fondamentale non solo per la teoria, ma anche per la pratica, poiché la natura di un accordo e la sua funzione possono cambiare drasticamente a seconda di quale intervallo si forma con la nota più grave (il basso).

\subsection{Dalla melodia all'armonia: la simultaneità dei suoni}

Finora abbiamo analizzato gli intervalli come relazioni tra coppie di suoni. Tuttavia, la musica occidentale, almeno a partire dal tardo Medioevo, non si limita a far suonare note una dopo l'altra (dimensione \textbf{orizzontale} o \textbf{melodica}), né a far suonare simultaneamente solo due note. La sua vera ricchezza risiede nella capacità di sovrapporre \textit{simultaneamente} tre o più suoni (dimensione \textbf{verticale} o \textbf{armonica}), creando strutture sonore complesse che hanno una loro logica interna e una loro funzione espressiva.

Questa sovrapposizione non è arbitraria o casuale. Come vedremo, esistono modi privilegiati di combinare i suoni che generano stabilità o tensione, che creano senso di movimento o di riposo. L'elemento base di questa costruzione verticale è l'\textbf{accordo}.


\begin{definizione}[Accordo]
	\label{def:accordo}
	Un \textbf{accordo} è la sovrapposizione simultanea di tre o più suoni. Nella pratica tonale, gli accordi si costruiscono prevalentemente sovrapponendo intervalli di terza. Un accordo di tre note è detto \textbf{triade}, uno di quattro note è detto \textbf{quadriade} (o accordo di settima), e così via.
\end{definizione}

La costruzione per terze sovrapposte non è casuale: come vedremo nel Capitolo \ref{chap:triadi}, essa riflette la struttura naturale della serie armonica e genera combinazioni di suoni che il nostro orecchio percepisce come unitarie e coerenti. Un accordo non è una semplice giustapposizione di note, ma un'entità sonora con una propria identità, funzione e tendenza risolutiva all'interno del sistema tonale.

\section{Il sistema tonale}
\label{sec:sistema_tonale}

Il sistema tonale non è una semplice successione di note, ma una \textbf{struttura gerarchica}. Ogni grado della scala ha un nome e una funzione specifica, creando una rete di tensioni e risoluzioni che guidano il discorso musicale.

\begin{definizione}[Gradi e funzioni della scala]
	\label{def:gradi_scala}
	I gradi della scala hanno nomi che ne descrivono la funzione (l'argomento sarà approfondito nel Capitolo \ref{chap:funzioni_armoniche}) e vengono denotati con numeri romani:
	\begin{itemize}
		\item \textbf{I grado (tonica):} il centro di gravità, il punto di massimo riposo.
		\item \textbf{V grado (dominante):} il polo di massima tensione, genera una forte aspettativa verso la tonica.
		\item \textbf{VII grado (sensibile):} così chiamata perché dista solo un semitono dalla tonica e "tende" a risolvere su di essa. È il motore principale della tensione della dominante.
		\item \textbf{IV grado (sottodominante):} crea una tensione più debole, "plagale", e spesso precede la dominante.
	\end{itemize}
	Gli altri gradi (II sopratonica, III mediante, VI sopradominante) hanno funzioni secondarie, spesso di preparazione a quelli principali.
\end{definizione}

La differenza principale tra il \textbf{modo maggiore} e il \textbf{modo minore} risiede nella diversa sequenza di toni e semitoni. In particolare, la distanza tra il I e il III grado (l'intervallo di terza) determina il "colore" del modo. Un'altra asimmetria cruciale riguarda la sensibile: nel modo maggiore è naturalmente presente, mentre nel modo minore naturale il VII grado è un tono sotto la tonica (si chiama infatti sottotonica). Per creare la necessaria tensione, questo grado viene quasi sempre innalzato cromaticamente, dando origine alla \textbf{scala minore armonica}.

 L'innalzamento di un semitono cromatico della sottotonica crea sì una sensibile, ma introduce anche un intervallo "problematico" di seconda aumentata tra sopradominante e sensibile che la melodia tende a evitare.

\section{Aritmetica e geometria tonale}
\label{sec:aritmetica_geometria}

Le relazioni tra le tonalità non sono casuali, ma seguono uno schema geometrico e matematico preciso, visualizzabile attraverso il \textbf{circolo delle quinte}.

	\begin{figure}[H]
	\centering
	% SCALA RIDOTTA per adattarsi alla pagina
	\begin{tikzpicture}[scale=3.6] 
		% --- Cerchi concentrici ---
		\draw[thick, gray!80] (0,0) circle (1.5cm);
		\draw[thick, gray!80] (0,0) circle (1.1cm);
		
		% --- Tonalità maggiori (cerchio esterno) ---
		\node at (90:1.3cm) {\textbf{Do}};
		\node at (60:1.3cm) {\textbf{Sol}};
		\node at (30:1.3cm) {\textbf{Re}};
		\node at (0:1.3cm) {\textbf{La}};
		\node at (330:1.3cm) {\textbf{Mi}};
		\node at (300:1.3cm) {\textbf{Si}};
		\node at (270:1.3cm) {\textbf{\Fadiesis /\Solbemolle }};
		\node at (240:1.3cm) {\textbf{\Rebemolle }};
		\node at (210:1.3cm) {\textbf{\Labemolle }};
		\node at (180:1.3cm) {\textbf{\Mibemolle }};
		\node at (150:1.3cm) {\textbf{\Sibemolle }};
		\node at (120:1.3cm) {\textbf{Fa}};
		
		% --- Tonalità minori (cerchio interno) ---
		\node[scale=0.85] at (90:0.8cm) {la m};
		\node[scale=0.85] at (60:0.8cm) {mi m};
		\node[scale=0.85] at (30:0.8cm) {si m};
		\node[scale=0.85] at (0:0.8cm) {fa\sharp\ m};
		\node[scale=0.85] at (330:0.8cm) {do\sharp\ m};
		\node[scale=0.85] at (300:0.8cm) {sol\sharp\ m};
		\node[scale=0.85] at (270:0.8cm) {re\sharp/mi\flat\ m};
		\node[scale=0.85] at (240:0.8cm) {si\flat\ m};
		\node[scale=0.85] at (210:0.8cm) {fa m};
		\node[scale=0.85] at (180:0.8cm) {do m};
		\node[scale=0.85] at (150:0.8cm) {sol m};
		\node[scale=0.85] at (120:0.8cm) {re m};
		
		% --- Elenco esplicito delle alterazioni (anello più esterno) ---
		% Raggi leggermente ridotti per un layout più compatto
		\node[gray, scale=0.7] at (90:1.75cm) {Nessuna};
		\node[gray, scale=0.7] at (60:1.75cm) {\Fadiesis };
		\node[gray, scale=0.7, align=center] at (30:1.75cm) {\Fadiesis , \Dodiesis };
		\node[gray, scale=0.7, align=center] at (0:1.75cm) {\Fadiesis , \Dodiesis , \Soldiesis };
		\node[gray, scale=0.7, align=center] at (330:1.75cm) {\Fadiesis , \Dodiesis , \Soldiesis , \Rediesis };
		\node[gray, scale=0.7, align=center] at (300:1.75cm) {\Fadiesis , \Dodiesis , \Soldiesis , \Rediesis , \Ladiesis };
		\node[gray, scale=0.65, align=center] at (270:1.85cm) {%
			$\substack{\text{\Fadiesis , \Dodiesis , \Soldiesis , \Rediesis , \Ladiesis , \Midiesis } \\ \text{\Sibemolle , \Mibemolle , \Labemolle , \Rebemolle , \Solbemolle , \Dobemolle }}$};
		\node[gray, scale=0.7, align=center] at (240:1.75cm) {\Sibemolle , \Mibemolle , \Labemolle , \Rebemolle , \Solbemolle };
		\node[gray, scale=0.7, align=center] at (210:1.75cm) {\Sibemolle , \Mibemolle , \Labemolle , \Rebemolle };
		\node[gray, scale=0.7, align=center] at (180:1.75cm) {\Sibemolle , \Mibemolle , \Labemolle };
		\node[gray, scale=0.7, align=center] at (150:1.75cm) {\Sibemolle , \Mibemolle };
		\node[gray, scale=0.7] at (120:1.75cm) {\Sibemolle };
		
		% --- Frecce indicatrici (posizionate più vicino al cerchio) ---
		\draw[->, very thick, red!70, line width=1.2pt] (100:2.1cm) arc (100:20:2.1cm);
		\node[red!70] at (60:2.35cm) {\textbf{QUINTE} $\rightarrow$};
		
		\draw[<-, very thick, blue!70, line width=1.2pt] (140:2.1cm) arc (140:220:2.1cm);
		\node[blue!70] at (180:2.35cm) {$\leftarrow$ \textbf{QUARTE}};
		
	\end{tikzpicture}
	\caption{Il circolo delle quinte e delle quarte con elenco esplicito delle alterazioni.}
	\label{fig:circolo_quinte_migliorato_compatto}
\end{figure}

Ogni passo sul circolo corrisponde a un salto di quinta giusta ascendente (o quarta giusta discendente). Questo schema non è solo una rappresentazione grafica, ma riflette le leggi fisiche e acustiche che abbiamo discusso: le tonalità più vicine sul circolo condividono più note e hanno più armonici in comune, e sono quindi percepite come più "vicine" anche dall'orecchio.

% ============================================================
% INTEGRAZIONE SCIENTIFICA - FISICA
% ============================================================
\begin{fisica}[La serie degli armonici come fondamento fisico della tonalità]
	\label{phys:serie_armonici_fondamento}
	
	Per comprendere perché un suono musicale non è una singola frequenza, dobbiamo partire dall'equazione delle \textbf{onde stazionarie}. Quando una corda vibrante o una colonna d'aria sono vincolate ai loro estremi, possono vibrare solo secondo particolari modi, chiamati \textbf{modi normali di vibrazione}.
	
	La condizione fisica per questi modi è che la lunghezza della corda $L$ deve contenere un numero intero di mezze lunghezze d'onda. Questo si esprime matematicamente come:
	\[
	L = n \cdot \frac{\lambda_n}{2} \quad \text{con } n = 1, 2, 3, \ldots
	\]
	
	Usando la relazione fondamentale tra velocità, frequenza e lunghezza d'onda ($v = \lambda \cdot f$), otteniamo le frequenze dei modi di vibrazione possibili:
	\[
	f_n = n \cdot f_1 = n \cdot \frac{v}{2L}
	\]
	
	dove $f_1$ è la \textbf{frequenza fondamentale} (primo armonico) e $f_n$ sono gli \textbf{armonici superiori}. Questa equazione ci dice che:
	\begin{itemize}
		\item Il secondo armonico ha frequenza $2f_1$ (un'ottava sopra la fondamentale)
		\item Il terzo armonico ha frequenza $3f_1$ (una quinta giusta sopra l'ottava)
		\item Il quarto armonico ha frequenza $4f_1$ (due ottave sopra)
		\item Il quinto armonico ha frequenza $5f_1$ (una terza maggiore sopra le due ottave)
	\end{itemize}
	
	\textbf{Rapporti di frequenza e intervalli musicali}
	
	Analizziamo ora i rapporti di frequenza tra questi armonici e la fondamentale:
	\begin{itemize}
		\item \textbf{Ottava}: rapporto $\frac{f_2}{f_1} = \frac{2:1}{1} = 2:1$
		\item \textbf{Quinta giusta}: rapporto $\frac{f_3}{f_2} = \frac{3f_1}{2f_1} = 3:2$
		\item \textbf{Terza maggiore}: rapporto $\frac{f_5}{f_4} = \frac{5f_1}{4f_1} = 5:4$
	\end{itemize}
	
	Questi rapporti semplici tra numeri interi piccoli (2:1, 3:2, 5:4) sono precisamente quelli che il nostro orecchio percepisce come i più \textbf{consonanti} e \textbf{stabili}. La fisica delle onde stazionarie giustifica quindi perché l'accordo maggiore (costruito con questi intervalli) sia la base naturale del sistema tonale: \textit{è insito nella natura stessa della vibrazione}.
	
	\textbf{Conseguenze per la teoria armonica}
	
	La serie armonica spiega perché:
	\begin{enumerate}
		\item La quinta giusta (3:2) appare prima della terza maggiore (5:4) nella serie: la quinta è un intervallo più "fondamentale"
		\item L'accordo maggiore (Do-Mi-Sol) ha rapporti approssimativamente 4:5:6, molto vicini a multipli consecutivi
		\item Gli intervalli con rapporti più complessi (come il tritono) sono percepiti come più dissonanti
	\end{enumerate}
\end{fisica}

% ============================================================
% INTEGRAZIONE SCIENTIFICA - MATEMATICA
% ============================================================
\begin{matematica}[L'aritmetica modulare come "orologio" del sistema tonale]
	\label{math:aritmetica_modulare}
	
	Il sistema tonale occidentale divide l'ottava in 12 semitoni, creando una struttura che si ripete identica a ogni ottava. Questa natura \textbf{ciclica} può essere descritta in modo elegante usando l'\textbf{aritmetica modulare}, lo stesso tipo di matematica che usiamo per leggere l'orologio.
	
	\textbf{Definizione dell'insieme $\mathbb{Z}_{12}$}
	
	Associamo a ogni nota una classe di altezza rappresentata da un numero da 0 a 11:
	\[
	\begin{array}{cccccccccccc}
	\text{Do} & \text{\Dodiesis } & \text{Re} & \text{\Rediesis } & \text{Mi} & \text{Fa} & \text{\Fadiesis } & \text{Sol} & \text{\Soldiesis } & \text{La} & \text{\Ladiesis } & \text{Si} \\
	0 & 1 & 2 & 3 & 4 & 5 & 6 & 7 & 8 & 9 & 10 & 11
	\end{array}
	\]
	
	L'insieme $\mathbb{Z}_{12} = \{0, 1, 2, \ldots, 11\}$ con l'\textbf{operazione di addizione modulo 12} forma quello che in matematica si chiama un \textbf{gruppo ciclico}.
	
	\textbf{Operazioni in $\mathbb{Z}_{12}$}
	
	L'addizione modulo 12 si comporta come l'aritmetica dell'orologio: quando superiamo 11, "torniamo indietro" a 0. Formalmente:
	\[
	a +_{12} b = (a + b) \mod 12
	\]
	
	\textbf{Esempio pratico}: Calcolo di intervalli
	
	\begin{itemize}
		\item \textbf{Quinta giusta ascendente}: un intervallo di quinta corrisponde a 7 semitoni.\\
		Da Do (0) a Sol: $0 + 7 = 7$ (Sol)\\
		Da Mi (4) a Si: $4 + 7 = 11$ (Si)\\
		Da Sol (7) a Re: $7 + 7 = 14 \equiv 2 \pmod{12}$ (Re)
		
		\item \textbf{Trasposizione}: Trasporre un pezzo di una terza maggiore (4 semitoni) significa sommare 4 a ogni nota.\\
		Accordo di Do maggiore: $(0, 4, 7)$ diventa $(4, 8, 11)$ = Mi maggiore
	\end{itemize}
	
	\textbf{Il circolo delle quinte come sequenza algoritmica}
	
	Il circolo delle quinte può essere generato da una semplice \textbf{funzione ricorsiva}:
	\[
	x_{n+1} = (x_n + 7) \mod 12
	\]
	
	Partendo da Do ($x_0 = 0$), generiamo:
	\begin{align*}
	x_1 &= (0 + 7) \mod 12 = 7 \quad \text{(Sol)} \\
	x_2 &= (7 + 7) \mod 12 = 2 \quad \text{(Re)} \\
	x_3 &= (2 + 7) \mod 12 = 9 \quad \text{(La)} \\
	x_4 &= (9 + 7) \mod 12 = 4 \quad \text{(Mi)} \\
	&\vdots
	\end{align*}
	
	Dopo 12 iterazioni, torniamo a 0 (Do), completando il circolo. Questo dimostra che il circolo delle quinte non è solo una convenzione teorica, ma una proprietà \textbf{matematica intrinseca} del sistema a 12 semitoni.
	
	\textbf{Proprietà interessanti}
	
	\begin{itemize}
		\item L'operazione "+7" genera tutti i 12 elementi di $\mathbb{Z}_{12}$ perché 7 e 12 sono \textbf{coprimi} ($\text{MCD}(7, 12) = 1$)
		\item Il circolo delle quarte si ottiene con l'operazione inversa: $x_{n+1} = (x_n + 5) \mod 12$, equivalente a $x_{n+1} = (x_n - 7) \mod 12$
		\item La simmetria enarmonica (\Fadiesis  = \Solbemolle ) corrisponde al fatto che $6 + 6 = 12 \equiv 0 \pmod{12}$: il tritono divide esattamente l'ottava a metà
	\end{itemize}
	
	\textbf{Applicazione pratica}
	
	Questa formalizzazione matematica non è solo un esercizio teorico: è la base degli \textbf{algoritmi} usati nei software di composizione musicale, nei traspositori automatici e negli strumenti che analizzano progressioni armoniche. Comprendere il sistema tonale come una struttura algebrica ci permette di manipolarlo in modo sistematico e preciso.
\end{matematica}
